\section{工件身份与复现}
\label{sec:artifact}
主数据集为 \path{frozen-trial/fitness.raw.json}，由归档提交 \code{584abd88}（tree \code{12586d3e}）对应的源码生成。表~\ref{tab:identity} 给出便于定位的短标识；证据清单提供完整哈希、大小和路径。源码归档中的全部 1003 个文件均已与导出版本核对。运行前后清单包含相同的 774 个路径及内容哈希，代码和设计修订均未变化。

\begin{table}[H]
\small\centering
\caption{主工件身份（SHA-256 前缀）。}
\label{tab:identity}
\begin{tabularx}{\columnwidth}{@{}Xl@{}}
\toprule
工件 & SHA-256 前缀\\
\midrule
原始试验报告 & \code{cf2575276eef583d}\\
保留的可执行文件，解压后 & \code{2ffaac58a3fe25ef}\\
可执行文件 gzip 归档 & \code{56b39d3c254318ad}\\
源码 tar 归档 & \code{1ed5157ce1165927}\\
语料定义 & \code{6cf6d8403c48ac25}\\
评价器定义 & \code{43e5455c5758bd11}\\
最终运行核验 & \code{e4c1887409b26ad6}\\
\bottomrule
\end{tabularx}
\end{table}

最初的冻结调用因缺少报告输出目录而以退出码 101 结束，没有产生主原始报告。该进程结束后，所需目录被创建；保留的可执行文件经恢复和哈希核对后直接运行，未重新编译。这次恢复运行以退出码 0 结束，并产生主报告。最初的失败调用仍按失败记录保留。最终核验器的 29 项检查全部通过，并明确处理了两种清单排序约定的差异；原始核验器与恢复执行回执仍可获取。

研究源码包包含汇总、逐场景与谓词分析 CSV、对应的来源哈希清单，以及重新生成矢量图的脚本。随附证据归档保留完整原始报告、源码 tar 归档、压缩可执行文件、核验器，以及各项独立的历史运行。重新分析使用这些保留数据，无需模型服务。单独开展的冻结契约重放为 22 个确切测试各保留一份回执，并核实其使用同一可执行文件且源码未变。历史观察不并入这些结果。

重新执行必须建立新的运行记录。保留的可执行文件面向 macOS arm64，并包含编译时绝对路径；在其他环境重新构建需要兼容的 Rust 工具链与依赖，也会产生新的可执行文件身份。工件文档提供命令与环境细节。公开存档和归档标识应对应实际发布的版本；本地复现包并不隐含这些事项已经完成。
